\documentclass[11pt,DIV=11]{scrartcl} % 
%Datum und Uhrzeit
\usepackage{datetime}

%Encoding
\usepackage[utf8]{inputenc}
\usepackage{lmodern}

%Graphikpakete
\usepackage{graphicx}
\usepackage[usenames,dvipsnames]{xcolor}
\usepackage{tikz}
\usetikzlibrary{arrows, snakes, backgrounds}
\usepackage{wrapfig}
\usepackage{dirtree}

%------------------------------------------------------------------
%% THI - Farbdefinitionen
%------------------------------------------------------------------

% Farben der THI
\definecolor{THIblue}{rgb}{0.0078,0.1176,0.4705}
\definecolor{AllianzBlue}{rgb}{0.07,0.20,0.53}
%%-------------Allgemeine Definitionen----------------------------------
% Farbige Aufwertung der berschriften
\addtokomafont{chapter}{\color{AllianzBlue}}
\addtokomafont{section}{\color{AllianzBlue}}
\addtokomafont{caption}{\color{AllianzBlue}}
\addtokomafont{subsection}{\color{AllianzBlue}}
\addtokomafont{subsubsection}{\color{AllianzBlue}}
\setkomafont{captionlabel}{\color{AllianzBlue}}


%Zitierumgebung
%\usepackage{cite}% Zitieren
\usepackage[ngerman]{babel}
\usepackage[backend=biber,]{biblatex}
%\usepackage{bibgerm}% Literatur in Deutscher DIN
\usepackage[babel,german=quotes]{csquotes}
\bibliography{ref/ref_liste} % Pfad und Datei der Ref-Datenbank

%URL-Umgebung
\usepackage{url}
\renewcommand{\UrlFont}{\small\color{THIblue}}

%Mathematik
\usepackage{amsmath}
\usepackage{amssymb}
\usepackage{mathtools}

%Tabellen
\usepackage{longtable}
\usepackage{multirow}

%Quellcode
\usepackage{listings}
\lstset{literate=%
    {Ö}{{\"O}}1
    {Ä}{{\"A}}1
    {Ü}{{\"U}}1
    {ß}{{\ss}}1
    {ü}{{\"u}}1
    {ä}{{\"a}}1
    {ö}{{\"o}}1
    {~}{{\textasciitilde}}1
}

\lstset{% general command to set parameter(s)
	basicstyle=\small\ttfamily,%\small, % print whole listing small
	keywordstyle=\color{THIblue}\bfseries,%\underbar,% underlined boldblack keywords
	identifierstyle=, % nothing happens
	commentstyle=\color{Aquamarine}, % white comments
	stringstyle=\color{Cerulean}\ttfamily, % typewriter type for strings
	showstringspaces=false, % no special string spaces
	frame=single,
	framexleftmargin=1.5em,
	xleftmargin=2em,
	numbers=left,
	numberstyle=\color{gray},
	numbersep=5pt,
	captionpos=b,
	breaklines=true}

\lstdefinelanguage{conf}
{
  % list of keywords
  basicstyle=\normalsize\ttfamily, %lineend
  sensitive=false, % keywords are not case-sensitive
  morecomment=[l]{\#}, % l is for line comment
  morecomment=[s]{/*}{*/}, % s is for start and end delimiter
  morestring=[b]" % defines that strings are enclosed in
}

\lstdefinelanguage{tikz}
{
  % list of keywords
  morekeywords={
	above,
	align,
  	begin,
	below,
  	caption,
	draw,
	end,
	every,
	fill,
	label,
  	line,
	minumum,
	node,
	rectangle,
	right,
	scale,
	size,
  	style,
	of,
  	width,
  	xshift,
  	yshift
  },
  sensitive=false, % keywords are not case-sensitive
  %alsodigit={0482},
  keywordstyle=\color{THIblue}\bfseries,
  morecomment=[l]{//}, % l is for line comment
  morecomment=[s]{/*}{*/}, % s is for start and end delimiter
  morestring=[b]" % defines that strings are enclosed in double quotes
}

%Microtypographie
\usepackage{microtype}

%Kopf und Fußzeilen
\usepackage{scrpage2} 	% Kopf & Fußzeile im KOMA Stil
\pagestyle{scrheadings}	% Aktiviert Verwendung vordefinierter Kolumnentitel
\clearscrheadfoot 		% alle Standard-Werte und Formatierungen löschen
\setkomafont{pagehead}{\scshape}	% Schriftart in Kopfzeile, \scshape = Kapitelchen
\automark[chapter]{section} % [linke Seite]{rechte Seite}
%\ohead{\def\pagestyle{PDTS}{\hrulefill\includegraphics[width = 6cm]{bilder/thi_logo_quer_cropped}}}
\ohead{\includegraphics[width = 6cm]{bilder/thi_logo_quer_cropped}}

%\ihead{\textsc{Abschlussarbeit}}
\ihead{\headmark}

%\setheadwidth[0pt]{textwithmarginpar}
\ofoot{\vspace{-0.3cm} \pagemark} 						
\ifoot{\vspace{-0.3cm} Dominik Gunther Florian Schlecht} 
				
%\setheadtopline{.4pt}				
\setheadsepline{.2pt}
\setfootsepline{.4pt}	% Trennlinie Fußzeile und Textkörper

%------------------------------------------------------------------
%% Längenanpassungen
%------------------------------------------------------------------
\setlength{\headsep}{10mm}		% Textabstand zur Kopfzeile
\setlength{\footskip}{15mm}		% Abstand zur Fußzeile
\setlength{\parindent}{0em}		% Einzug nach Absatz

%%-------------Allgemeine Definitionen----------------------------------
% Farbige Aufwertung der berschriften
\addtokomafont{chapter}{\color{THIblue}}
\addtokomafont{section}{\color{THIblue}}
\addtokomafont{caption}{\color{THIblue}}
\addtokomafont{subsection}{\color{THIblue}}
\addtokomafont{subsubsection}{\color{THIblue}}
\setkomafont{captionlabel}{\color{THIblue}}

%Hyperref
\usepackage[
		pdftex,
		linkcolor=THIblue,			% Farbe der Verlinkung
		linktocpage=true,			% Im TOC wird Seitenzahl verlinkt(true),bzw. Text(false)
		colorlinks=true,			% 'true' keinen Kasten um Link
		citecolor=THIblue,
%		pdfhighlight=/P,
%		bookmarks,
%		hyperfigures=true,
%		citebordercolor={0 0 1},
%		linkbordercolor={0 0 1},
%		menubordercolor={0 0 1},
%		backref=true,
%		pagebackref=true,
%		bookmarksopen,
%		bookmarksnumbered,
%		pdfpagelabels=false,
%		pdfstartpage=1,
%		pdfstartview=Fit,			% Modus beim Öffnen (Fit = An Seitengröße anpassen)
		pdftitle={Umgehen von USB-Deskriptor basierten USB-Policies am Beispiel einer virtuellen Umgebung},
		pdfauthor={Dominik Schlecht},
%		pdfstartview=Fit,
%		pdfdisplaydoctitle=true,
%		plainpages=false
			]{hyperref}  

%------------------------------------------------------------------
% Wichtige Definition für Aufteilung von Formelverzeichnis und Abkürzungsverzeichnis
%------------------------------------------------------------------

% Nomenklaturverzeichnis, Formelzeichenliste Anpassungen für nomenclature: damit lassen sich zwei getrennte Symbolverzeichnisse anlegen, ziehmlich cool!
\usepackage[intoc,compatible,german]{nomencl}	
		\renewcommand{\nomgroup}[1]{	\ifthenelse{\equal{#1}{A}}{\item[{\normalfont\sffamily\bfseries\LARGE\textcolor[rgb]{0,0.112,0.47}{Abkürzungen{\phantom{\Huge $\frac{\frac{\frac{A}{a}}{a}}{\frac{a}{a}}$}}}}]}{	\ifthenelse{\equal{#1}{A}}{\item[{\normalfont\sffamily\bfseries\LARGE\textcolor[rgb]{0,0.112,0.47}{Formelzeichen{\phantom{\Huge $\frac{A}{\frac{a}{a}}$}}}}]}{}}}
		
%------------------------------------------------------------------
%% Anpassung von Abständen, Längen vom Nomenclaturverzeichnis (Abkürzungs- und Formelverzeichnis)
%------------------------------------------------------------------

% Abstand zwischen Einträgen im Symbolverzeichnis (-\parsep = 0)
\setlength{\nomitemsep}{-\parsep} 
\setlength{\nomlabelwidth}{5em}	
\renewcommand{\nomlabel}[1]{#1 \dotfill}  % Punkte im zwischen Nummer und Kapiteleintrag

% Sub-Figures mit Caption
\usepackage{subcaption}
\newcommand{\frage}[1]{\makebox[\textwidth]{%
\renewcommand{\arraystretch}{2.0}
\begin{tabularx}{\textwidth}{|X|}
  \hline
  #1 \\
  \hline
  \line(1,0){420}\\
  \hline \hline
\end{tabularx}%
\renewcommand{\arraystretch}{1.0}
}}

\newcommand{\frageJaNein}[1]{\makebox[\textwidth]{%
\renewcommand{\arraystretch}{2.0}
\begin{tabularx}{\textwidth}{|X|S|S|}
  \hline
  #1 & Ja & Nein\\
  \hline
  \line(1,0){350} & $\square$ & $\square$	\\
  \hline \hline
\end{tabularx}%
\renewcommand{\arraystretch}{1.0}
}}

\newcommand{\frageOpt}[4]{\makebox[\textwidth]{%
\renewcommand{\arraystretch}{2.0}
\begin{tabularx}{\textwidth}{|X|X|X|}
  \hline
  \multicolumn{3}{|c|}{#1}\\
  \hline
  #2 $\square$ & #3 $\square$ & #4 $\square$\\
  \hline \hline
\end{tabularx}%
\renewcommand{\arraystretch}{1.0}
}}

\newcommand{\frageJaNeinKurz}[1]{\makebox[\textwidth]{%
\renewcommand{\arraystretch}{1.5}
\begin{tabularx}{\textwidth}{|X|S|S|}
  \hline
  #1 & Ja $\square$ & Nein $\square$\\
  \hline \hline
\end{tabularx}%
\renewcommand{\arraystretch}{1.0}
}}
\usepackage[parfill]{parskip}
\begin{document}
\section{Allgemeines}
\subsection{Ansprechpartner}
%%allg_anspr_web_app

\subsection{Eingrenzung}
Welche Art/en von Pentest/s sollen durchgeführt werden?
\begin{itemize}
	\item[%%allg_pen_art_wapt] Web-Application
	\item[%%allg_pen_art_npt] Network
	\item[%%allg_pen_art_se] Social Engineering
	\item[%%allg_pen_art_wl] Wireless
	\item[%%allg_pen_art_phys] Physical
\end{itemize}

\subsection{Allgemeine Fragen}\label{ref:allg_frag}
\makebox[\textwidth]{%
\renewcommand{\arraystretch}{2.0}
\begin{tabularx}{\textwidth}{|X|S|S|}
  \hline
  Ist der Test für eine spezielle Compliance-Anforderung notwendig? & Ja & Nein\\
  \hline
  %%alg_compliance	\\
  \hline \hline
\end{tabularx}%
\renewcommand{\arraystretch}{1.0}
}\\

\makebox[\textwidth]{%
\renewcommand{\arraystretch}{2.0}
\begin{tabularx}{\textwidth}{|X|}
  \hline
  Wann soll der Test statt finden? \\
  \hline
  %%allg_termin \\
  \hline \hline
\end{tabularx}%
\renewcommand{\arraystretch}{1.0}
}

\makebox[\textwidth]{%
\renewcommand{\arraystretch}{2.0}
\begin{tabularx}{\textwidth}{|X|X|X|}
  \hline
  \multicolumn{3}{|c|}{In welchen Zeiträumen soll der Test durchgeführt werden?}\\
  \hline
  %%alg_zeitraum \\
  \hline \hline
\end{tabularx}%
\renewcommand{\arraystretch}{1.0}
}
\newpage

\section{Web Application Penetration Test}\label{ref:WebAppPenTest}

\makebox[\textwidth]{%
\renewcommand{\arraystretch}{2.0}
\begin{tabularx}{\textwidth}{|X|S|S|}
  \hline
  Wird der Quellcode der Applikation/Webseite zugänglich gemacht? & Ja & Nein\\
  \hline
  %%wapt_quell_zug	\\
  \hline \hline
\end{tabularx}%
\renewcommand{\arraystretch}{1.0}
}

\makebox[\textwidth]{%
\renewcommand{\arraystretch}{2.0}
\begin{tabularx}{\textwidth}{|X|}
  \hline
  Wie viele Web-Applikationen sind In-Scope? \\
  \hline
  %%wapt_anz_web_app \\
  \hline \hline
\end{tabularx}%
\renewcommand{\arraystretch}{1.0}
}

\makebox[\textwidth]{%
\renewcommand{\arraystretch}{2.0}
\begin{tabularx}{\textwidth}{|X|}
  \hline
  Wie viele Login-Systeme sind In-Scope? \\
  \hline
  %%wapt_anz_login_sys \\
  \hline \hline
\end{tabularx}%
\renewcommand{\arraystretch}{1.0}
}

\makebox[\textwidth]{%
\renewcommand{\arraystretch}{2.0}
\begin{tabularx}{\textwidth}{|X|}
  \hline
  Wie viele statische Seiten sind ca. In-Scope? \\
  \hline
  %%wapt_anz_stat_seiten \\
  \hline \hline
\end{tabularx}%
\renewcommand{\arraystretch}{1.0}
}

\makebox[\textwidth]{%
\renewcommand{\arraystretch}{2.0}
\begin{tabularx}{\textwidth}{|X|}
  \hline
  Wie viele dynamische Seiten sind ca. In-Scope? \\
  \hline
  %%wapt_anz_dyn_seiten \\
  \hline \hline
\end{tabularx}%
\renewcommand{\arraystretch}{1.0}
}

\makebox[\textwidth]{%
\renewcommand{\arraystretch}{2.0}
\begin{tabularx}{\textwidth}{|X|S|S|}
  \hline
  Soll Fuzzing gegen die Applikation/en eingesetzt werden? & Ja & Nein\\
  \hline
  %%wapt_fuzzing	\\
  \hline \hline
\end{tabularx}%
\renewcommand{\arraystretch}{1.0}
}

\makebox[\textwidth]{%
\renewcommand{\arraystretch}{2.0}
\begin{tabularx}{\textwidth}{|X|S|S|}
  \hline
  Soll der Penetrations-Test aus verschiedenen Rollen durchgeführt werden? & Ja & Nein\\
  \hline
  %%wapt_vers_rollen	\\
  \hline \hline
\end{tabularx}%
\renewcommand{\arraystretch}{1.0}
}

\makebox[\textwidth]{%
\renewcommand{\arraystretch}{2.0}
\begin{tabularx}{\textwidth}{|X|S|S|}
  \hline
  Sollen Password-Scans auf die Webseite durchgeführt werden? & Ja & Nein\\
  \hline
  %%wapt_pwd_scan	\\
  \hline \hline
\end{tabularx}%
\renewcommand{\arraystretch}{1.0}
}

\section{Network Penetration Test}\label{ref:NetPenTest}

\makebox[\textwidth]{%
\renewcommand{\arraystretch}{2.0}
\begin{tabularx}{\textwidth}{|X|}
  \hline
  Was ist das Ziel des Penetrations-Test? \\
  \hline
  %%npt_ziel \\
  \hline \hline
\end{tabularx}%
\renewcommand{\arraystretch}{1.0}
}

\makebox[\textwidth]{%
\renewcommand{\arraystretch}{2.0}
\begin{tabularx}{\textwidth}{|X|}
  \hline
  Wie viele IP-Adressen sollen getestet werden? \\
  \hline
  %%npt_anz_ips \\
  \hline \hline
\end{tabularx}%
\renewcommand{\arraystretch}{1.0}
}

\makebox[\textwidth]{%
\renewcommand{\arraystretch}{2.0}
\begin{tabularx}{\textwidth}{|X|S|S|}
  \hline
  Sind Techniken im Einsatz, die die Resultate verfälschen könnten? (WAF, IPS etc.?) & Ja & Nein\\
  \hline
  %%npt_verteidigung	\\
  \hline \hline
\end{tabularx}%
\renewcommand{\arraystretch}{1.0}
}

\makebox[\textwidth]{%
\renewcommand{\arraystretch}{2.0}
\begin{tabularx}{\textwidth}{|X|}
  \hline
  Wie ist das Vorgehen bei einem gelungenen Angriff? \\
  \hline
  %%npt_gel_ang \\
  \hline \hline
\end{tabularx}%
\renewcommand{\arraystretch}{1.0}
}

\makebox[\textwidth]{%
\renewcommand{\arraystretch}{2.0}
\begin{tabularx}{\textwidth}{|X|S|S|}
  \hline
  Soll versucht werden lokale Admin-Rechte zu erlangen und tiefer in das Netz vorzudringen? & Ja & Nein\\
  \hline
  %%npt_nutz_lok_admin	\\
  \hline \hline
\end{tabularx}%
\renewcommand{\arraystretch}{1.0}
}

\makebox[\textwidth]{%
\renewcommand{\arraystretch}{2.0}
\begin{tabularx}{\textwidth}{|X|S|S|}
  \hline
  Sollen Angriffe auf gefundene Passwort-Hashes durchgeführt werden? & Ja & Nein\\
  \hline
  %%npt_pwd_hashes	\\
  \hline \hline
\end{tabularx}%
\renewcommand{\arraystretch}{1.0}
}

\section{Social Engineering}\label{ref:SocEngi}

\makebox[\textwidth]{%
\renewcommand{\arraystretch}{2.0}
\begin{tabularx}{\textwidth}{|X|S|S|}
  \hline
  Gibt es eine vollständige Liste von E-Mail-Adressen, die für den Test verwendet werden können? & Ja & Nein\\
  \hline
  %%se_emails	\\
  \hline \hline
\end{tabularx}%
\renewcommand{\arraystretch}{1.0}
}

\makebox[\textwidth]{%
\renewcommand{\arraystretch}{2.0}
\begin{tabularx}{\textwidth}{|X|S|S|}
  \hline
  Gibt es eine vollständige Liste von Telefon-Nummern, die für den Test verwendet werden können? & Ja & Nein\\
  \hline
  %%se_tels	\\
  \hline \hline
\end{tabularx}%
\renewcommand{\arraystretch}{1.0}
}

\makebox[\textwidth]{%
\renewcommand{\arraystretch}{2.0}
\begin{tabularx}{\textwidth}{|X|S|S|}
  \hline
  Ist das Einsetzen von Social Engineering zum Überwinden physikalischer Sicherheitseinrichtungen erlaubt? & Ja & Nein\\
  \hline
  %%se_phys	\\
  \hline \hline
\end{tabularx}%
\renewcommand{\arraystretch}{1.0}
}

\makebox[\textwidth]{%
\renewcommand{\arraystretch}{2.0}
\begin{tabularx}{\textwidth}{|X|}
  \hline
  Wie viele Personen sollen ca. getestet werden? \\
  \hline
  %%se_anz_pers \\
  \hline \hline
\end{tabularx}%
\renewcommand{\arraystretch}{1.0}
}

\section{Wireless Network Penetration Test}\label{ref:WirNetPen}

\makebox[\textwidth]{%
\renewcommand{\arraystretch}{2.0}
\begin{tabularx}{\textwidth}{|X|}
  \hline
  Wieviele Funk-Netzwerke sind im Einsatz? \\
  \hline
  %%wl_anz_fnw \\
  \hline \hline
\end{tabularx}%
\renewcommand{\arraystretch}{1.0}
}

\makebox[\textwidth]{%
\renewcommand{\arraystretch}{2.0}
\begin{tabularx}{\textwidth}{|X|S|S|}
  \hline
  Gibt es eine Gäste WLAN? Wenn ja, wie ist dieses Umgesetzt? & Ja & Nein\\
  \hline
  %%wl_gastwlan	\\
  \hline \hline
\end{tabularx}%
\renewcommand{\arraystretch}{1.0}
}

\makebox[\textwidth]{%
\renewcommand{\arraystretch}{2.0}
\begin{tabularx}{\textwidth}{|X|}
  \hline
  Welche Verschlüsselung wird für die Netzwerke genutzt? \\
  \hline
  %%wl_encr \\
  \hline \hline
\end{tabularx}%
\renewcommand{\arraystretch}{1.0}
}

\makebox[\textwidth]{%
\renewcommand{\arraystretch}{2.0}
\begin{tabularx}{\textwidth}{|X|S|S|}
  \hline
  Sollen nicht-firmen-Geräte im WLAN aufgespürt werden? & Ja & Nein\\
  \hline
  %%wl_fremgeraet	\\
  \hline \hline
\end{tabularx}%
\renewcommand{\arraystretch}{1.0}
}

\makebox[\textwidth]{%
\renewcommand{\arraystretch}{2.0}
\begin{tabularx}{\textwidth}{|X|S|S|}
  \hline
  Soll Netz-Attacken gegen Clients durchgeführt werden? & Ja & Nein\\
  \hline
  %%wl_att_clients	\\
  \hline \hline
\end{tabularx}%
\renewcommand{\arraystretch}{1.0}
}

\makebox[\textwidth]{%
\renewcommand{\arraystretch}{2.0}
\begin{tabularx}{\textwidth}{|X|}
  \hline
  Wie viele Clients nutzen das WLAN ca.? \\
  \hline
  %%wl_clients \\
  \hline \hline
\end{tabularx}%
\renewcommand{\arraystretch}{1.0}
}

\section{Physical Penetration Test}\label{ref:PhysPen}

\makebox[\textwidth]{%
\renewcommand{\arraystretch}{2.0}
\begin{tabularx}{\textwidth}{|X|}
  \hline
  Wie viele Einrichtungen sollen getestet werden? \\
  \hline
  %%phys_anz_einr \\
  \hline \hline
\end{tabularx}%
\renewcommand{\arraystretch}{1.0}
}

\makebox[\textwidth]{%
\renewcommand{\arraystretch}{2.0}
\begin{tabularx}{\textwidth}{|X|S|S|}
  \hline
  Wird die Einrichtung mit anderen Parteien geteilt? & Ja & Nein\\
  \hline
  %%phys_anz_part	\\
  \hline \hline
\end{tabularx}%
\renewcommand{\arraystretch}{1.0}
}

\makebox[\textwidth]{%
\renewcommand{\arraystretch}{2.0}
\begin{tabularx}{\textwidth}{|X|S|S|}
  \hline
  Muss Sicherheitspersonal umgangen werden? & Ja & Nein\\
  \hline
  %%phys_sich_pers \\
  \hline \hline
\end{tabularx}%
\renewcommand{\arraystretch}{1.0}
}

\makebox[\textwidth]{%
\renewcommand{\arraystretch}{2.0}
\begin{tabularx}{\textwidth}{|X|S|S|}
  \hline
  Wird das Sicherheitspersonal durch einen Dritten gestellt? & Ja & Nein\\
  \hline
  %%phys_sich_pers_dritt	\\
  \hline \hline
\end{tabularx}%
\renewcommand{\arraystretch}{1.0}
}

\makebox[\textwidth]{%
\renewcommand{\arraystretch}{2.0}
\begin{tabularx}{\textwidth}{|X|S|S|}
  \hline
  Ist das Sicherheitspersonal bewaffnet? & Ja & Nein\\
  \hline
  %%phys_sich_pers_waffen	\\
  \hline \hline
\end{tabularx}%
\renewcommand{\arraystretch}{1.0}
}

\makebox[\textwidth]{%
\renewcommand{\arraystretch}{2.0}
\begin{tabularx}{\textwidth}{|X|S|S|}
  \hline
  Ist der Einsatz von körperlicher Gewalt durch das Sicherheitspersonal gestattet? & Ja & Nein\\
  \hline
  %%phys_sich_pers_gewalt	\\
  \hline \hline
\end{tabularx}%
\renewcommand{\arraystretch}{1.0}
}

\makebox[\textwidth]{%
\renewcommand{\arraystretch}{2.0}
\begin{tabularx}{\textwidth}{|X|}
  \hline
  Wie viele Eingänge gibt es zu der/den Einrichtung/en? \\
  \hline
  %%phys_anz_eing \\
  \hline \hline
\end{tabularx}%
\renewcommand{\arraystretch}{1.0}
}

\makebox[\textwidth]{%
\renewcommand{\arraystretch}{2.0}
\begin{tabularx}{\textwidth}{|X|S|S|}
  \hline
  Ist das knacken von Schlössern oder fälschen von Schlüsseln erlaubt? & Ja & Nein\\
  \hline
  %%phys_knacken	\\
  \hline \hline
\end{tabularx}%
\renewcommand{\arraystretch}{1.0}
}

\makebox[\textwidth]{%
\renewcommand{\arraystretch}{2.0}
\begin{tabularx}{\textwidth}{|X|}
  \hline
  Wie groß ist die Fläche ungefähr? \\
  \hline
  %%phys_flaeche \\
  \hline \hline
\end{tabularx}%
\renewcommand{\arraystretch}{1.0}
}

\makebox[\textwidth]{%
\renewcommand{\arraystretch}{2.0}
\begin{tabularx}{\textwidth}{|X|S|S|}
  \hline
  Sind alle physikalischen Sicherheitsmaßnahmen dokumentiert und werden zur Verfügung gestellt? & Ja & Nein\\
  \hline
  %%phys_whitebox	\\
  \hline \hline
\end{tabularx}%
\renewcommand{\arraystretch}{1.0}
}

\makebox[\textwidth]{%
\renewcommand{\arraystretch}{2.0}
\begin{tabularx}{\textwidth}{|X|S|S|}
  \hline
  Werden Video-Kameras verwendet? & Ja & Nein\\
  \hline
  %%phys_cam	\\
  \hline \hline
\end{tabularx}%
\renewcommand{\arraystretch}{1.0}
}

\makebox[\textwidth]{%
\renewcommand{\arraystretch}{2.0}
\begin{tabularx}{\textwidth}{|X|S|S|}
  \hline
  Werden diese Kameras durch Dritte verwaltet? & Ja & Nein\\
  \hline
  %%phys_cam_dritte	\\
  \hline \hline
\end{tabularx}%
\renewcommand{\arraystretch}{1.0}
}

\makebox[\textwidth]{%
\renewcommand{\arraystretch}{2.0}
\begin{tabularx}{\textwidth}{|X|S|S|}
  \hline
  Soll versucht werden, die aufgezeichneten Daten zu löschen? & Ja & Nein\\
  \hline
  %%phys_loeschen	\\
  \hline \hline
\end{tabularx}%
\renewcommand{\arraystretch}{1.0}
}

\makebox[\textwidth]{%
\renewcommand{\arraystretch}{2.0}
\begin{tabularx}{\textwidth}{|X|S|S|}
  \hline
  Gibt es ein Alarm-System? & Ja & Nein\\
  \hline
  %%phys_alarm	\\
  \hline \hline
\end{tabularx}%
\renewcommand{\arraystretch}{1.0}
}

\makebox[\textwidth]{%
\renewcommand{\arraystretch}{2.0}
\begin{tabularx}{\textwidth}{|X|S|S|}
  \hline
  Gibt es einen Stillen Alarm? & Ja & Nein\\
  \hline
  %%phys_alarm_still	\\
  \hline \hline
\end{tabularx}%
\renewcommand{\arraystretch}{1.0}
}

\makebox[\textwidth]{%
\renewcommand{\arraystretch}{2.0}
\begin{tabularx}{\textwidth}{|X|}
  \hline
  Welche Ereignisse lösen den Alarm aus? \\
  \hline
  %%phys_alarm_ausloeser \\
  \hline \hline
\end{tabularx}%
\renewcommand{\arraystretch}{1.0}
}

\section{Questions for Systems Administrators}\label{ref:QuesSysAdmin}

\makebox[\textwidth]{%
\renewcommand{\arraystretch}{2.0}
\begin{tabularx}{\textwidth}{|X|S|S|}
  \hline
  Gibt es Systeme, die als instabil angesehen werden (alte Patch-Stände, Legacy Systeme etc.)? & Ja & Nein\\
  \hline
  %%sysadm_altsysteme	\\
  \hline \hline
\end{tabularx}%
\renewcommand{\arraystretch}{1.0}
}

\makebox[\textwidth]{%
\renewcommand{\arraystretch}{2.0}
\begin{tabularx}{\textwidth}{|X|S|S|}
  \hline
  Gibt es Systeme von Dritten, die ausgeschlossen werden müssen oder für die weitere Genehmigungen notwendig sind? & Ja & Nein\\
  \hline
  %%sysadm_drittsysteme	\\
  \hline \hline
\end{tabularx}%
\renewcommand{\arraystretch}{1.0}
}

\makebox[\textwidth]{%
\renewcommand{\arraystretch}{2.0}
\begin{tabularx}{\textwidth}{|X|}
  \hline
  Was ist die Durchschnittszeit zur Wiederherstellung der Funktionalität eines Services? \\
  \hline
  %%sysadm_wiederherstellungszeit \\
  \hline \hline
\end{tabularx}%
\renewcommand{\arraystretch}{1.0}
}

\makebox[\textwidth]{%
\renewcommand{\arraystretch}{2.0}
\begin{tabularx}{\textwidth}{|X|S|S|}
  \hline
  Ist eine Monitoring-Software im Einsatz? & Ja & Nein\\
  \hline
  %%sysadm_monitoring	\\
  \hline \hline
\end{tabularx}%
\renewcommand{\arraystretch}{1.0}
}

\makebox[\textwidth]{%
\renewcommand{\arraystretch}{2.0}
\begin{tabularx}{\textwidth}{|X|}
  \hline
  Welche sind die kritischsten Applikationen? \\
  \hline
  %%sysadm_krit_appl \\
  \hline \hline
\end{tabularx}%
\renewcommand{\arraystretch}{1.0}
}

\makebox[\textwidth]{%
\renewcommand{\arraystretch}{2.0}
\begin{tabularx}{\textwidth}{|X|S|S|}
  \hline
  Werden in einem regelmäßigen Turnus Backups erstellt und getestet? & Ja & Nein\\
  \hline
  %%sysadm_backup	\\
  \hline \hline
\end{tabularx}%
\renewcommand{\arraystretch}{1.0}
}

\section{Questions for Business Unit Managers}\label{ref:QuesBUM}

\makebox[\textwidth]{%
\renewcommand{\arraystretch}{2.0}
\begin{tabularx}{\textwidth}{|X|S|S|}
  \hline
  Ist die Führungsebene über den Test informiert? & Ja & Nein\\
  \hline
  %%bum_mgmt_inform	\\
  \hline \hline
\end{tabularx}%
\renewcommand{\arraystretch}{1.0}
}

\makebox[\textwidth]{%
\renewcommand{\arraystretch}{2.0}
\begin{tabularx}{\textwidth}{|X|}
  \hline
  Welche Daten stellen das größte Risiko dar, falls diese manipuliert werden? \\
  \hline
  %%bum_krit_daten \\
  \hline \hline
\end{tabularx}%
\renewcommand{\arraystretch}{1.0}
}

\makebox[\textwidth]{%
\renewcommand{\arraystretch}{2.0}
\begin{tabularx}{\textwidth}{|X|S|S|}
  \hline
  Gibt es Testfälle, die die Funktionalität der Services prüfen und belegen können? & Ja & Nein\\
  \hline
  %%bum_monitoring	\\
  \hline \hline
\end{tabularx}%
\renewcommand{\arraystretch}{1.0}
}

\makebox[\textwidth]{%
\renewcommand{\arraystretch}{2.0}
\begin{tabularx}{\textwidth}{|X|S|S|}
  \hline
  Sind "Disaster Recovery Procedures" vorhanden? & Ja & Nein\\
  \hline
  %%bum_disaster_recovery	\\
  \hline \hline
\end{tabularx}%
\renewcommand{\arraystretch}{1.0}
}

\end{document}